\section{Background}
\label{sec:ch2_bg}

\myx{It seems that much of this subsection overlaps with the prior one. Maybe you can discuss more about the unique challenges of applying caching in agentic setting?}

The key-value (KV) cache in transformer models stores intermediate attention states of processed tokens so that each new token need not recompute attention over the entire sequence. However, this cache normally resets between separate API calls. To reuse computation across calls, engines like vLLM employ batching and paged attention, treating the KV cache as pageable memory \cite{kwon2023vllm}. SGLang \cite{zheng2024sglangefficientexecutionstructured} introduces RadixAttention, storing caches in a radix tree keyed by prefixes and scheduling requests to maximize cache hits.

Yet these optimizations target exact prefix reuse. PromptCache  generalizes reuse by letting developers annotate prompt modules, but it cannot adapt to content or order changes \cite{gim2024promptcachemodularattention}. CacheBlend advances the state-of-the-art by making cache reuse position-independent. It treats the prompt as composed of chunks, checks for each chunk if it has stored KV state, and when reassembling chunks recomputes only a selection of HKVD (High-KV-Deviation) tokens to update cross-attention and positional encodings, restoring the impact of cross-attention on the forward-attention matrix \cite{yao2025cacheblendfastlargelanguage}. The observation in which HKVD tokens tend to remain deviated through different layers also made it possible for efficient selection and KV recomputation of such tokens. This allows seamless merging of caches for arbitrarily ordered prompt segments, resulting in near-perfect cache utilization and dramatic latency savings with negligible accuracy loss.
